\documentclass[12pt]{exam}
\usepackage[T1]{fontenc}
\usepackage[utf8]{inputenc}
\usepackage[lithuanian]{babel}
\usepackage{graphicx}
\usepackage[a4paper, left=1cm, right=1cm, top=2.5cm, bottom=1cm]{geometry}
\usepackage{amsmath}

% Antraštė
\pagestyle{headandfoot}
\runningheadrule
\firstpageheader{Vardas: \rule{4cm}{0.4pt}}{Kontrolinis darbas: \\Išvestinių taikymai}{Fibonačio variantas}
\runningheader{Meilė išgelbės pasaulį}{}{}
\firstpagefooter{}{}{}
\runningfooter{}{}{}
\pointsdroppedatright          % rodyti taškus dešinėje paraštėje
\bracketedpoints               % [5] vietoj 5  (nebūtina)
\setlength{\rightpointsmargin}{1.6cm} % kiek į dešinę nustumti taškus
\renewcommand{\points}{taškai} % daugiskaita
\renewcommand{\pointname}{Taškai:} % kai nenaudojama paraštė
\marginpointname{\ taškai}     % tekstas šalia skaičiaus paraštėje: [5 taškai]
\renewcommand{\dottedlinefillheight}{0.7cm}
\renewcommand{\points}{taškai}
\renewcommand{\pointname}{Taškai:}
\newcommand{\atsakymas}{%
    \par\vspace{0.5cm}
    \textbf{Atsakymas:} \framebox[5cm]{\rule{0pt}{0.7cm}}
}
\begin{document}

\begin{questions}
\question[6]
Consider the expansion of $\left(3x^{2}-\dfrac{k}{x}\right)^{9}$, where $k>0$.
The coefficient of the term in $x^{6}$ is $6048$. Find the value of $k$.
\vspace{5cm}
\atsakymas



\question Raskite funkcijų išvestines:
\begin{parts}
\part[2] $f(x) = \sqrt{3-5x ^{2}}+ \cos(\ln x)- \frac{1}{e^{4+x- x ^{3} } }$     \droppoints
\part[2] $g(x) = \frac{\cos 2x \cos 6x- \sin 6x \sin 2x }{ 2\sin 4x \cos 4x} $      \droppoints
\end{parts}

\vspace{5cm}
\atsakymas

\question
Remdamiesi funkcijos $f$ išvestinės grafiku nustatykite:
\begin{parts}
    \part[2] funkcijos $ f$ monotoniškumo intervalus.
        \droppoints
    \part[2] funkcijos $f$ minimumo ir maksimumo taškus.
    \droppoints
\end{parts} 

\includegraphics[scale = 0.3]{image2.png} % čia tavo paveikslėlis

\vspace{5cm}
\atsakymas




\question Nustatykite funkcijos $f(x) = \frac{3x}{1+x ^{2} }  $ 
\begin{parts}
    \part[4] ekstremumo taškus ir ekstremumus.     \droppoints
    \part[2] didžiausią ir mažiausią reikšmes intervale $x \in [0;5] $.     \droppoints 
\end{parts}
\vspace{5cm}
\atsakymas

\question[5] Tiesė liečia funkcijos $f(x) = x(x ^{2} -3)+1 $ grafiką taške, kurio abscisė lygi $0$ . 
\begin{parts}
    \part[2] Parašykite šios tiesės lygtį.    \droppoints
    \part[2] Dešimtųjų tikslumu nustatykite kampo dydį, kurį sudaro tiesė su ordinačių ašimi.    \droppoints
    \part[3] Apskaičiuokite plotą figūros, kurią koordinačių plokštumoje apriboja ši tiesė ir koordinačių\\ ašys.    
    \droppoints
\end{parts}

\vspace{5cm}
\atsakymas

\question[5] Reikia pagaminti kūgio formos indą. Kūgio sudaromoji turi būti lygi $ 20 \sqrt{3 }\;cm$ ir į indą turi tilpti didžiausias galimas skysčio kiekis. Kokio ilgio turi būti šio kūgio aukštinė?
\droppoints
\vspace{5cm}
\atsakymas


\end{questions}

\end{document}
